% !TEX root = chapter-standalone.tex

\section{Powerset}
\SYNDEF{powerset}
\label{sec:power-set}

\begin{ctdefinition}[Power set]
    \label{def:power-set}
    Given a set~\setA, we can form a new set whose elements are precisely all the subsets of~\setA.
    This new set is called the \emph{powerset} of~\setA; we denote it by~$\powerset \setA$.
\end{ctdefinition}

For example, if~$\setA = \makeset{ \sapple, \scarrot, \sgrapes }$, then its \SY{powerset} is
\begin{equation}
    \powerset \setA = \makeset{ \Emptyset, \makeset{ \sapple }, \makeset{ \scarrot }, \makeset{ \sgrapes}, \makeset{ \sapple, \scarrot }, \makeset{ \scarrot, \sgrapes }, \makeset{ \sapple, \sgrapes}, \makeset{ \sapple, \scarrot, \sgrapes } }.
\end{equation}

\begin{exercise}
    Can you count how many elements the \SY{powerset}~$\powerset \setA$ has in the following cases?
    \begin{enumerate}
        \item $\setA = \makeset{ \sapple } $.
        \item $\setA = \makeset{ \sapple, \scarrot } $.
        \item $\setA = \makeset{ \sapple, \scarrot, \sgrapes } $.
        \item $\setA = \Emptyset $.
    \end{enumerate}
    Can you guess a general formula for the size of the \SY{powerset} of a finite set?
\end{exercise}

\begin{solution}
    We have:
    \begin{enumerate}
        \item 2.
        \item 4.
        \item 8.
        \item 1.
    \end{enumerate}
    In general, the size of~$\powerset \setA$ is~$2$ to the power of the size of~\setA.
\end{solution}

Now suppose we fix a set~\setA for a moment.
Given~$\subA \setin\powerset \setA$, the \emph{complement} of~$\subA$ with respect to~\setA is
\begin{equation}
    \setA \setcomplement \subA = \makeset{ \ela \setin\setA \mid \ela \notsetin \subA },
\end{equation}
which is again an element of~$\powerset \setA$.
In situations where it is evident which ambient set~\setA we are working with, the notation~$\compl{\subA}$ is sometimes used instead of~$\setA \setcomplement \subA$.

We also note that the operations of union and intersection, when restricted to~$\powerset \setA$, again produce elements of~$\powerset \setA$.
That is, if~$\subA, \subB \setin\powerset \setA$, then~$\subA \setunion \subB \setin\powerset \setA$, and similarly~$\subA \setintersection \subB \setin\powerset \setA$.

The operations~$\setunion$,~$\setintersection$, and~$\compl{( - )}$ on~$\powerset \setA$ obey various rules which are useful to be familiar with.
For example,
\begin{equation}
    \compl{(\subA \setintersection \subB)} = (\compl{\subA}) \setunion (\compl{\subB}).
\end{equation}
Can you state more such rules?
A useful visual aid for such calculations are so-called Venn diagrams (\cref{fig:venn-union}, \cref{fig:venn-intersection}, \cref{fig:venn-complement}).

\begin{marginfigure}
    \centering
    \includesag{venn-union}
    \caption{Venn diagram for union operation.}
    \label{fig:venn-union}
\end{marginfigure}

\begin{marginfigure}
    \centering
    \includesag{venn-intersection}
    \caption{Venn diagram for intersection operation.}
    \label{fig:venn-intersection}
\end{marginfigure}

\begin{marginfigure}
    \centering
    \includesag{venn-complement}
    \caption{Venn diagram for complement operation.}
    \label{fig:venn-complement}
\end{marginfigure}

\begin{gradedexercise}[\exname{DistributingSubsets}]
    \label{ex:distributing-subsets}
    Let~\setA be a set, and let~$\subA, \subB, \subC \setsubseteq \setA$ be subsets.
    Prove that
    \begin{equation}
        \subA \setintersection (\subB \setunion \subC) = (\subA \setintersection \subB) \setunion (\subA \setintersection \subC).
    \end{equation}
\end{gradedexercise}

\solutionof{DistributingSubsets}




\todotext{2026-07-20, JL: i've now, for the moment, moved the sections on "union and intersection" to be here 'subordinate' to the powerset construction, but originally they were mean to discuss the more general operations of union and intersection that -- i think -- are allowed in ZFC foundations between any two sets; from the type theory point of view this is not allowed though, i think, and is perhaps directly counter to the whole idea of sets as separate types. Need to research a bit to see how arbitrary unions/intersection do or do not fit into the dependent type theory framework.}

\subsection{Union and intersection}
The union of two sets is the set containing precisely those elements which come from either of the two.
\begin{ctdefinition}[Union of sets]
    \label{def:union-of-sets}
    \SYNDEF{union of sets}
    Given sets~\setA and~\setB, their \emph{union} is a new set, denoted~$\setA \setunion \setB$, characterized by
    \begin{equation}
        \prfdoubleperiod{
            \ela \setin(\setA \setunion \setB)
        }{
            \ela \setin\setA \quad \boolor \quad \ela \setin\setB
        }
    \end{equation}
\end{ctdefinition}

For example, if~$\setA = \makeset{ \sapple, \sgrapes }$ and~$\setB = \makeset{ \stea }$, then
\begin{equation}
    \setA \setunion \setB = \makeset{ \sapple, \sgrapes, \stea }.
\end{equation}

% Or, if~$\setA = \makeset{ \sapple, \sgrapes, \scarrot }$ and~$\setB = \makeset{ \sapple, \scarrot, \swater }$, then
% \begin{equation}
%     \setA \setintersection \setB = \makeset{ \sapple, \scarrot }.
% \end{equation}

The intersection of two sets is the set of elements common to both.

\begin{ctdefinition}[Intersection of sets]
    \label{def:intersection-of-sets}
    \SYNDEF{intersection of sets}
    The \emph{intersection} of sets~\setA and~\setB, denoted~$\setA \setintersection \setB$, is the set characterized by
    \begin{equation}
        \prfdoubleperiod{
            \ela \setin(\setA \setintersection \setB)
        }{
            \ela \setin\setA \quad \booland \quad \ela \setin\setB
        }
    \end{equation}
\end{ctdefinition}

For example, if~$\setA = \makeset{ \sapple, \sgrapes, \scarrot, \stea }$ and~$\setB = \makeset{ \stea, \sgrapes, \swater }$, then
\begin{equation}
    \setA \setintersection \setB = \makeset{ \stea, \sgrapes }.
\end{equation}

\begin{exercise}
    Prove that union and intersection of sets are associative operations.
\end{exercise}

\begin{solution}
    We start by the union operation $\setunion$.
    We need to prove
    \begin{equation*}
        \setA \setunion (\setB \setunion \setC)= (\setA \setunion \setB) \setunion \setC.
    \end{equation*}
    We have:
    \begin{equation*}
        \begin{split}
                    & \ela \setin \setA \setunion (\setB \setunion \setC)                     \\
            \Eqv \  & \ela \setin \setA \boolor (\ela \setin \setB \boolor \ela \setin \setC) \\
            \Eqv \  & (\ela \setin \setA \boolor \ela \setin \setB)\boolor \ela \setin \setC  \\
            \Eqv \  & \ela \setin (\setA \setunion \setB) \setunion \setC.
        \end{split}
    \end{equation*}
    We now continue with the intersection operation $\setintersection$.
    We need to prove
    \begin{equation*}
        \setA \setintersection (\setB \setintersection \setC)= (\setA \setintersection \setB) \setintersection \setC.
    \end{equation*}
    We have:
    \begin{equation*}
        \begin{split}
                    & \ela \setin \setA \setintersection (\setB \setintersection \setC)         \\
            \Eqv \  & \ela \setin \setA \booland (\ela \setin \setB \booland \ela \setin \setC) \\
            \Eqv \  & (\ela \setin \setA \booland \ela \setin \setB)\booland \ela \setin \setC  \\
            \Eqv \  & \ela \setin (\setA \setintersection \setB) \setintersection \setC.
        \end{split}
    \end{equation*}
    Essentially, we have used the associativity of the $\booland$ and $\boolor$ connectives.
\end{solution}

\begin{exercise}
    Prove that union and intersection of sets are commutative operations.
\end{exercise}

\begin{solution}
    We start by the union operation $\setunion$.
    We need to prove that $\setA \setunion \setB = \setB \setunion \setA$.
    We have:
    \begin{equation*}
        \begin{split}
                    & \ela \setin \setA \setunion \setB           \\
            \Eqv \  & \ela \setin \setA \boolor \ela \setin \setB \\
            \Eqv \  & \ela \setin \setB \boolor \ela \setin \setA \\
            \Eqv \  & \ela \setin \setB \setunion \setA.
        \end{split}
    \end{equation*}
    We continue with the intersection operation $\setintersection$.
    We need to prove that $\setA \setintersection \setB = \setB \setintersection \setA$.
    We have:
    \begin{equation*}
        \begin{split}
                    & \ela \setin \setA \setintersection \setB     \\
            \Eqv \  & \ela \setin \setA \booland \ela \setin \setB \\
            \Eqv \  & \ela \setin \setB \booland \ela \setin \setA \\
            \Eqv \  & \ela \setin \setB \setintersection \setA.
        \end{split}
    \end{equation*}
    Essentially, we have used the commutativity of the $\booland$ and $\boolor$ connectives.
\end{solution}

The definitions of union and intersection above are for two sets, $\setA$ and $\setB$.
We can also build the union or intersection of any finite number of sets, or even an infinite collection of sets.
Let's look at the finite case first.

Given sets $\setAn{1}, \setAn{2}, \dots, \setAn{n}$, their union $\bigsetunion_{i \in \makeset{1,\dots, n}} \setAn{i}$ is defined by
\begin{equation}
    \prfdoubleperiod{ \ela \setin \bigsetunion_{i \in \makeset{1,\dots, n}} \setAn{i} }{\exists i \in \makeset{1,\dots, n} \colon \ela \setin\setAn{i}}
\end{equation}
Alternative notations for this union are also $\bigsetunion_{i = 1}^{n} \setAn{i}$ and $\bigsetunion \makeset{ \setAn{i} \mid i \in \makeset{1,\dots, n}}$.
The latter notation lends itself well for the generalization of the union operation to any arbitrary collection of sets.
If $\setofsetsA$ is a collection of sets (it might have two elements (say, sets~\setA and~\setB), or it might have 108 elements, or it might have infinitely many elements) we define $\bigsetunion \setofsetsA$ by
%
\begin{equation}
    \prfdoubleperiod{ \ela \setin\bigsetunion \setofsetsA }{\exists \setA \setin\setofsetsA \colon \ela \setin\setA}
\end{equation}
%
This notation for the union of an arbitrary collection of sets is related to our previous definition for two sets~\setA and~\setB via
\begin{equation}
    \setA \setunion \setB = \bigsetunion \makeset{ \setA, \setB }.
\end{equation}

\begin{remark}
    A slightly confusing thing might be the following; it has to do with how we use variables and the fact that for sets, elements can appear at most once.
    If we write for example ``$\setAn{1}, \setAn{2}, \dots, \setAn{n}$'', then we are a priori syntactically speaking about $n$ distinct sets.
    However, this notation does not exclude the possibility that perhaps $\setAn{1} = \setAn{2} = \dots = \setAn{n} = \setA$, for instance.
    In this special concrete case, the set of sets $\makeset{ \setAn{i} \mid i \in \makeset{1,\dots, n}}$ will have only one single element, namely $\setA$, even though, in terms of notation, it might look like there are more elements.
\end{remark}

\begin{remark}
    When we write ``$\setAn{1}, \setAn{2}, \dots, \setAn{n}$'', we are using the numbers $1, \dots, n$ to \emph{index} these $n$ (not necessarily non-equal) sets.
    Indexing means giving them distinct names.
    We don't necessarily need to use natural numbers to index a collection of sets (or any other things, for that matter) -- we can use any other set as an index! The main point is that the index set (let's call it $\stylesets{I}$) should have precisely as many elements as we wish to have distinct ``names'' for the sets we are indexing.
    Let's look at some examples.

    For instance, for a collection of four sets, we might name them $\setAn{1}$, $\setAn{2}$, $\setAn{3}$, $\setAn{4}$ using the index set $\stylesets{I} = \makeset{1,2, 3, 4}$, but we also might name them $\setAn{north}$, $\setAn{south}$, $\setAn{east}$, $\setAn{west}$ using the index set $\stylesets{I} = \makeset{north, south, east, west}$.
    With such small index sets, it is easy list all the sets involved; however in general the notation $\makeset{\setAn{i}}_{i \in \stylesets{I}}$ is used.
    For example, if $\stylesets{I} = \wnumbers$ or $\stylesets{I} = \reals$, we denote the respective corresponding indexed collections of sets by $\makeset{\setAn{k}}_{k \in \wnumbers}$ or $\makeset{\setAn{\lambda}}_{\lambda \in \reals}$, for example.
    An indexed collection of sets is also sometimes called a family of sets.
\end{remark}

Similar to how the operation of union may be generalized to any arbitrary collection of sets, so too the operation of intersection.
Given a collection $\setofsetsA$ of sets, we define the intersection~$\bigsetintersection \setofsetsA$ by
%
\begin{equation}
    \prfdoubleperiod{\ela \setin\bigsetintersection \setofsetsA}{\forall \setA \setin\setofsetsA \colon \ela \setin\setA}
\end{equation}
This notation for the intersection of an arbitrary collection of sets is related to our previous definition for two sets~\setA and~\setB via
\begin{equation} \setA \setintersection \setB = \bigsetintersection \makeset{ \setA, \setB }.
\end{equation}

\todojira{609}{\bernina: include some sort of a non-arbitrary (and useful??) example to make more concrete and illustrate?}

